\cotefile{49}
\scanpage{2}
\uline{Picard des \unsure{schèmas|schémas} \unsure{abèliens|abéliens}}
% archive sheet number: 1

\scanpage{3}
\uline{Théorème} Soit $A$ un schéma abélien projectif sur $S$. Alors $\underline{\mathrm{Pic}}^{\tau}\deleted{\unsure{A|X}}_{A/S}\inserted{=\underline{\mathrm{Pic}}^{0}_{A/S}}$ est un schéma abélien sur $S$, ouvert et fermé dans $\underline{\mathrm{Pic}}_{A/S}$ ayant $=$ dimension relative sur $S$ que $A$.
% slanted
\margin{Si on ne veut pas se servir de la nullité de la torsion, on trouve que $\underline{\mathrm{Pic}}_{X/S}$ est simple sur $S$ en les pts de $\underline{\mathrm{Pic}}^{0}_{A/S}$, \deleted{\unsure{et on}} \inserted{et on} voit alors que $\varphi$ est un morphisme \uline{plat}, donc ouvert, donc $\underline{\mathrm{Pic}}^{0}_{A/S}$ est aussi ouvert dans $\underline{\mathrm{Pic}}_{A/S}$, c'est donc un schéma abélien sur $S$ par le stratagème \unsure{indiqué|ind.} \unsure{ci-dessous|ci-dessus} \ldots.}

Comme $\deleted{\unsure{X}}A$ est simple sur $S$ et que ses fibres \uline{ont des \unsure{Ner. Sév.|Nér. Sév.|Ner. Sev.} sans torsion}, on sait que $\underline{\mathrm{Pic}}^{\tau}_{\deleted{\unsure{X}}A/S}$ est propre sur $S$, et fermé dans $\underline{\mathrm{Pic}}_{A/S}$. Il est en tous cas \uline{ouvert} dans $\underline{\mathrm{Pic}}^{\tau}_{A/S}$.

Il \unsure{reste|veut} à montrer qu'il est \uline{simple} sur $S$ (on sait déjà que ses fibres sont universellement connexes), et a la bonne dimension.

Soit donc $\underline L$ un Module ample sur $A$ rel. $S$, il définit une section $s$ de $\underline{\mathrm{Pic}}_{A/S}$ sur $S$, et un \inserted{morphisme}
\deleted{\unsure{définit|désigne} par translation un sous-schéma ouvert et fermé}
\[
\deleted{P=s+\underline{\mathrm{Pic}}^{0}_{A/S}}
\]
\deleted{de $\underline{\mathrm{Pic}}_{A/S}$. Or on a un morphisme}
\[
\varphi_s:\ A\longrightarrow \deleted{\unsure{P}}A'\deleted{\unsure{0|o}}=\underline{\mathrm{Pic}}^\tau_{X/S}
\]
obtenu posant
\[
\varphi_s(t)=\text{classe de }\tau_{-t}^{*}(\underline L)\,\underline L^{-1}
\]
[Comme $\varphi_s$ transforme section unité en section unité, c'est un hom de schémas en groupes (\unsure{réf|rif})]. \deleted{\unsure{Donc}} L'hom $A_0\to A'_0$ défini par $\varphi$ est \uline{surjectif} [car le noyau est fini d'après le critère d'amplitude de Weil, d'autre part
\[
\dim A'_0\underset{\text{\unsure{déformation du fibre|définition du fibre|déformation du Pic|déformation de Pic}}}{\underbrace{\leq}} \dim H^1(A_0,\underline O_{A_0})
\underset{\substack{\text{th. de Hopf }+\\ \text{nullité de }H^i(A_0,\underline O_{A_0})\text{ pour }i>\dim A_0}}{\underbrace{\leq}}\dim A_0
\]
% archive sheet number: 2

\scanpage{4}
et comme $\dim A_0\leq\dim A'_0$ par ce qui précède, on a égalité partout ; \unsure{ceci|ce} prouve en particulier que \uline{$A'_0$ est simple} ], et aussi
\[
\dim H^{1}(A_0,\underline O_{A_0})=n,\quad\text{d'où la structure totale de }H^{\unsure{x|*}}(A_0,\underline O_{A_0})\unsure{\big]|.}
\]

Comme on vient de voir que $A'_0$ est simple, on voit que pour tt $x'\in A'_0$, $\exists\,x\in A_0$ au dessus et on aura nécessairement $\underline O_{\deleted{\unsure{X}}A'_0,x'}\unsure{\text{---}|\longrightarrow}\underline O_{A_0,x}$ \uline{injectif}.

Or on a le

\uline{Lemme} $X$ et $Y$ de type fini sur $S$ loc. noeth. \inserted{\unsure{avec|soient} $F,G$ \unsure{quacoh.|quacoh|quasi-coh.} sur $X,Y$} $\varphi:X\to Y$ un $S$-morphisme, $s\in S$, \inserted{$G\to F$ un $\varphi$-morphisme}, \deleted{\unsure{on suppose}} $y\in Y$ au dessus de $y$, on suppose a) $\deleted{\unsure{G}} \unsure{F|G}$ est plat sur $\unsure{S_y|S,|S}$ b/ $G$ coh.

% slanted
\margin{mettre aussi le cas où l'hypo. noeth. \deleted{\unsure{de finitude|de Noether}} \unsure{peut-être|peut être} sera remplacée par l'hypothèse que $Y$ \deleted{\unsure{réduit}} est local à radical nilpotent \ldots}

b) Soit $\varphi_s:X_s\to Y_s$ induit par $\varphi$, \deleted{\unsure{alors}} tout \unsure{élémt|élém$^{\mathrm t}$|élément} de $G_s$ dont l'image inverse en les $x\in\varphi_s^{-1}(y)$ s'annule est \inserted{nulle} \deleted{$G_s\to\varphi_{s,*}(F_s)$ est injectif en $y$.} \uline{Alors}

$G$ est plat sur $S$ en $y$.

On utilise le critère par les gradués associés.

On applique ceci au $\varphi:A\to A_{\deleted{\unsure{0}}}'$ et un pt quelconque de $A'$, $F=\underline O_A$, $G=\underline O_{A'}$.

Cela prouve que $A'$ \unsure{est|$\to$} plat sur $S$ \uline{cqfd}.
% archive sheet number: 3

\scanpage{5}
\noindent\llap{\smash{\rule[-1.8\baselineskip]{0.4pt}{1.7\baselineskip}}\quad}\uline{Corollaire 1} \deleted{$R^{\unsure{1|\prime|i}}f_*(\underline O_A)$ est \unsure{l}} Le faisceau $\underline O_A$ est coh. plat pour le \deleted{\unsure{$f$|$\varphi$}} morphisme de projection $A\to S$. \deleted{L'\unsure{Algèbre}} On a un isomorphisme d'Algèbres
\[
R^{*}f_*(\underline O_A)\simeq\bigwedge\nolimits^{*} R^{1}f_*(\underline O_A)
\]
\noindent\llap{\smash{\rule[-1.2\baselineskip]{0.4pt}{1.5\baselineskip}}\quad}où $R^1f_{\unsure{*|\,}}(\underline O_A)$ est loc. libre sur $S$, de rang égal à la dimension relative de $A$ sur $S$.

\uline{Dém.} On sait déjà que l'énoncé est vrai quand $S$ est le spectre d'un corps (cf ci-dessus). \deleted{\unsure{Prouvons}} D'autre part, il est connu que le fait que $\underline{\mathrm{Pic}}_{A/S}$ soit simple sur $S$ le long de la section unité $\varepsilon$ implique que
\[
R^1f_*(\underline O_A)=\underline{\mathrm{Hom}}(\deleted{\unsure{\varepsilon^*(\Omega^1)|\Omega^1_{S}|\varepsilon^*(\Omega^1_{A/S})}}\,\underline{\Omega}^{\unsure{1|\prime}}_\varepsilon,\underline O_S)
\]
est localement libre sur $S$, et de façon plus précise que $\underline O_A$ est coh. plat pour $f$ en dim 1.

Considérons \unsure{l'e|le} \deleted{\unsure{hom}} \inserted{diagramme} canonique
\diagram{\begin{tikzcd}[ampersand replacement=\&,column sep=large]
\bigwedge^{*}R^1f_*(\underline O_A)_y \arrow[r,"\alpha"] \arrow[d,"\mathrm{surj}"] \& R^{\unsure{*|x}}f_*(\underline O_A)_y \arrow[d] \\
\bigwedge^{*}H^1(X_y,\underline O_{X_y}) \arrow[r,"{\substack{\deleted{\mathrm{surj}}\\[-1pt]\sim}}"] \& H^{*}(X_y,\underline O_{X_y})
\end{tikzcd}}
\deleted{\unsure{la flèche qui est|le carré qui est}}
\[
\deleted{R^{*}f_*(\underline O_A)_y\longrightarrow{}}
\]
cela montre que $R^*f_*(\underline O_A)_y\to H^{*}(X_y,\underline O_{X_y})$ est \uline{surjectif}, d'où le fait que $\underline O_A$ est coh. plat pour $f$. Comme $\alpha$ est un hom de Modules loc. libres qui induit un isom sur les fibres
% archive sheet number: 4

\scanpage{6}
réduites, c'est un isom, cqfd.

\uline{Remarque} J'ignore si le corollaire reste valable sans supposer $A$ projectif sur $S$, tt comme j'ignore si $\underline{\mathrm{Pic}}_{A/S}$ existe dans ce cas, et s'il est \unsure{encore|ou non|même} simple sur $S$ le long de la section unité\ldots.

\uline{Corollaire 2} Si $\underline L$ est ample rel. \unsure{à|:} $S$, alors (et alors \unsure{seult|seult.}) $\varphi_{\underline L}:A\unsure{\text{---}|\longrightarrow} A'$ est un morphisme surjectif, \unsure{et même|étant|et} \inserted{\unsure{fini}} fidèlement plat ; son noyau est un groupe fini et plat sur $\unsure{S|s}$.
% slanted
\margin{énoncer sous forme de condition \unsure{néc.|nèc.} et suffisante sur $\underline L$.}

\uline{Corollaire} 3 $\underline{\mathrm{Pic}}^+_{\deleted{\unsure{X}}A/S}$ est ouvert \uline{et fermé}

[le ouvert est toujours le cas ; le fermé est spécial aux schémas abéliens].

En effet, si $\varphi_{\underline L}$ induit un hom surjectif en $s$, il induit un hom surjectif aux pts de $\overline s$.

\uline{Corollaire} 4 Existence de $\underline{\text{Nér}}^{\unsure{\varepsilon|\xi|S|\text{Sév}}}_{\deleted{\unsure{X}}A/S}$
% archive sheet number: 5

\scanpage{7}
Nous allons étudier de plus près le morphisme
\[
\underline{\mathrm{Pic}}_{A/S}\longrightarrow\underline{\mathrm{Hom}}_{S\text{-gr}}(A,A')\tag*{$(*)$}
\]
\margin{hom pointé, \deleted{\unsure{ou|de|$\leftarrow$}} hom de groupes}
i.e.
\[
\underline{\mathrm{Pic}}_{A/S}\times_S A\longrightarrow A'\tag*{$(**)$}
\]
défini par
\[
(\xi\unsure{,|;}t)\rightsquigarrow t_*(\xi)-\xi
\]

\noindent\llap{\smash{\rule[-3.6ex]{0.4pt}{4ex}}\quad}\uline{Théorème} L'hom $(*)$ est un hom de $S$-Groupes \unsure{?|$\xi$|:} (i.e. $(**)$ est un accouplement), et son noyau est \uline{exactement} $A'$, \deleted{\unsure{la}} \inserted{et même :} \deleted{\unsure{dès que|pour que} \unsure{l'on a|une}} une section $\xi$ de $\underline{\mathrm{Pic}}_{A/S}$ sur $S$ est une \deleted{\unsure{est nulle|distincte} \unsure{la|de} \unsure{loi}} section de $A'$ \unsure{ssi|ses} on a
\[
\boxed{\pi^*(\xi)=\mathrm{pr}_1^*(\xi)+\mathrm{pr}_2^*(\xi)}
\]
où $\pi:A\times_{\unsure{S|s}} A\longrightarrow A$ est la \unsure{loi de} composition.

\noindent\llap{\smash{\rule[-1.8ex]{0.4pt}{2ex}}\quad}\uline{Corollaire} On déduit de $(*)$ par passage au quotient une \uline{immersion fermée}
\[
\underline{\text{\unsure{NérSév|NérSé}}}_{A/S}\longrightarrow\underline{\mathrm{Hom}}_S(A,A')\tag*{$(**)$}
\]
% slanted
\margin{$\underline{\text{NérSév}}_{A/S}=\underline{\mathrm{Pic}}_{A/S}/\underline{\mathrm{Pic}}^0_{A/S}$ par définition}
compatible avec les structures de groupe.

Démonstration du théorème. L'additivité en $t$ est immédiate grâce au fait que le morphisme $A\unsure{\text{---}|\longrightarrow}A'$ défini par $s$ transforme section unité en section unité. L'additivité en $s$ est de toute façon triviale. Le morphisme $A'\times A\longrightarrow A'$ induit par $(**)$
% archive sheet number: 6

\scanpage{8}
induit des morphismes \inserted{sur} $A'\times_S e$ et $e\times_S A$ qui sont triviaux. Donc il est trivial sur \deleted{\unsure{$S$|$\xi$}} $A\times A'$, \deleted{\unsure{d'après}} étant \unsure{multiplicatif} sur $A\times A'$, qui est un schéma abélien.
% slanted
\margin{ici on utilise le fait que $A'$ est \uline{abélien}}
Donc on est ramené à prouver que le morphisme $(**)$ \deleted{\unsure{est une}} a un noyau nul, et il suffit d'ailleurs de prouver que \deleted{\unsure{il}} c'est une immersion fermée. Comme $\underline{\text{\unsure{NérSév|NérSé}}}_{A/S}$ \deleted{est} donc \deleted{\unsure{tout}} \unsure{son -préschéma|son sous-préschéma} est non ramifié sur $S$, on est ramené à prouver \deleted{que les} pour le premier pt que les fibres du noyau de $\underline{\text{NérSév}}_{A/S}\to\underline{\mathrm{Hom}}_{S\text{-gr}}(A,A')$ sont ensemblistement réduites à un pt. Donc on est ramené au cas où $S$ est le spectre d'un corps alg. clos, et à prouver ceci :

\noindent\llap{\smash{\rule[-3.7ex]{0.4pt}{4.2ex}}\quad}\uline{Lemme} $A$ var. \unsure{abélienne|abélien} sur un corps $k$ \unsure{alg. clos|alg clos} $\xi$ une classe de diviseurs sur $A$, \unsure{invariante|invariant} par translation. Alors $\xi$ est alg. équivalente à $0$

C'est connu, et sera prouvé plus bas \ldots..
% archive sheet number: 7

\scanpage{9}
% slanted
\margin{L'énoncer \unsure{indifféremt|indifféremment} avec $\xi$ ou $L$}
\deleted{\unsure{$\underline L$ est|Lemme}} On va prouver ici que si $\xi\in\mathrm{Pic}(A/S)$, alors $\varphi_\xi=0$ \unsure{sss|ssi} on a
\[
\pi^*(\xi)-\mathrm{pr}_1^*(\xi)-\mathrm{pr}_2^*(\xi)=0
\]
En effet, considérons le diagramme
\diagram{\begin{tikzcd}[ampersand replacement=\&]
 A \& A\times A \arrow[l,"\mathrm{pr}_2"'] \arrow[r,"\pi"] \arrow[d,"\mathrm{pr}_1"] \& A\\
 \& A \&
\end{tikzcd}}
Dire que $\varphi_\xi=0$ signifie aussi que, après le chang$^{\mathrm{mt}}$ de base \deleted{\unsure{induit}} \unsure{sur|de|si} $S_1=A\unsure{\text{---}|\longrightarrow} S$, posant $A_1=A\times_S S_1$, et considérant la section universelle $t_1$ de $A_1$ sur $S_1$, on aura $t_1^*(\xi_1)=\xi_1$ (où $\xi_1=\xi\times_S S_1$).
% square at left of the lines beginning “posant”
\diagram{\begin{tikzcd}[ampersand replacement=\&]
A \arrow[d] \& A^{\unsure{\prime|\,}}_1=A\times_{\unsure{1|S}} A \arrow[l,"\mathrm{pr}_2"'] \arrow[d,"\mathrm{pr}_1"] \\
S \& A=S_1 \arrow[l,"\unsure{q|g}"]
\end{tikzcd}}
Or \deleted{$\unsure{t_1^*(\xi_1)}\,\unsure{t_1^*}(\xi_1)\unsure{=}\,t_1^*(\xi_1)={}$} Si $\xi$ est défini par $\underline L$ sur $A/S$, \inserted{trivial sur la section unité.} alors $\xi_1$ par $\mathrm{pr}_2^*(L)$, alors $t_1^*(\xi_1)$ est défini par
\[
\begin{aligned}
t_1^*(\mathrm{pr}_2^*(L))&=(\mathrm{pr}_2\,\mathrm{transl}(t_1))^*(\underline L)\\
&=\pi^*(\underline L).
\end{aligned}
\]
On doit donc \deleted{\unsure{prouver}} écrire que $\pi^*(\underline L)$ et $\mathrm{pr}_2^*(\underline L)$ définissent \unsure{la|le} $=$ \unsure{classe|élt} de $\mathrm{Pic}(A\times_S A/A)$, ou ce qui revient au même que
% slanted
\margin{l'écrire \unsure{tensoriellement}}
$\underbrace{\pi^*(\underline L)-\mathrm{pr}_1^*(\underline L)-\mathrm{pr}_2^*(\underline L)}$ définit \unsure{une|un} \unsure{classe|élt} \uline{\unsure{nulle|nul}} de $\mathrm{Pic}(A\times_S A/A)$. Comme ce faisceau induit un faisceau trivial sur la section unité de $A\times_S A$ sur $A$ :
\[
i_1^*(\text{---})=(\underbrace{\pi i_1}_{\mathrm{id}})^*(\underline L)-(\underbrace{\mathrm{pr}_1 i_1}_{\mathrm{id}})^*(\underline L)-(\underbrace{\mathrm{pr}_2 i_1}_{\varepsilon_{A/S}})^*(\underline L)=0
\]
% archive sheet number: 8

\scanpage{10}
\unsure{cela|ceci|cette} signifie que le faisceau en question est trivial.
% Faint sideways show-through beneath the writing.
% archive sheet number: 9

